\honest{Terminal Values and Instrumental Values}{Eliezer Yudkowsky}{2007}

You want chocolate from the Publix, so you pick up the car keys. Then the radio says an earthquake has destroyed all the chocolate there, and every step loses its point at once. People rarely lose track of plans like this. Humans are experts at planning, not experts on planning. I have ``noticed'' that when they talk about goals in the abstract, they confuse means and ends, or, ``more formally'', ``instrumental values'' and ``terminal values''. \nb{The distinction is old. Philosophers call it instrumental versus final, or intrinsic, value, and trace the idea of something ``good for its own sake'' to Aristotle. The post does not claim otherwise.}

Part of the trouble, it seems to me, is that English uses the one word ``want'' for saving my sister's life and for giving her penicillin. My first definitions are in English. Instrumental values are wanted only for their consequences: I want penicillin for my sister because I expect it to cure her, and if I expected it to melt her like the Wicked Witch I would fight to keep it from her. Terminal values are wanted regardless of consequences. But if saving my sister would get the Earth swallowed by a black hole, I would not save her, so even her life looks conditional on consequences.

``So forget English.'' I set out an ``ideal Bayesian decision system'': outcomes, actions, a utility for each outcome, and a probability of each outcome given each action. With utility 1 if my sister lives and 0 if she dies, penicillin has an expected utility of 0.9 and no penicillin 0.3, so the system chooses penicillin. A plan of several steps can be handled ``by letting each Action stand for a whole sequence'', at the cost of an exponentially large space, which also loses structure a human planner uses: if one first act is shooting your own foot off, a person discards every sequence that starts that way. So there are ``a few minor complications,'' but it is surprisingly useful to consider the absurdly simple version first. \nb{A fixed sequence is enough only when nothing learned along the way could change the later steps. In the chocolate example the radio news arrives partway through. Each action must then be a strategy, a rule for what to do after each possible observation, and there are more of those still.}

Now the use of the model. A philosopher says that everyone is ultimately selfish: a mother who helps her son only wants to believe he is well. You ask about a mother who dies pushing her son out of the path of a truck. The philosopher answers that she did it because ``she valued'' that choice, for ``the feeling of importance she attached to that decision.'' The right answer is ``TYPE ERROR: No constructor found for Expected\_Utility -> Utility.'' Expected utility belongs to actions and utility to outcomes; both map to real numbers, as wind speed and temperature do, but that does not make them the same.

If all your utilities were over your own states of mind, you would steer the future toward your happiness and be indifferent between any two futures with the same state of mind, and you would rarely die for another. When the philosopher's answer shifts from utilities of outcomes to the value of the decision, it makes a jump that sounds the same in English but would be an error in a typed language. That a decider chose what it rated highest does not ``say anything whatsoever about where it steers the future.'' \nb{This is a standard objection to psychological egoism: if self-interest means satisfying all one's preferences, egoism becomes trivially true. But the type error fits only one reading of the philosopher. On another, the feeling of importance is a state of the mother's mind at the moment of choosing, and so part of an outcome, which the model allows: she preferred dying with that feeling to living without it. That is the egoist's usual reply. What counts against it is evidence about motives, such as Daniel Batson's experiments, and the post gives none.}

To save your son you must imagine him saved, but what you imagine is your son, not your own imagining. The imagination is a quotation, like ``snow'' against snow. To steer toward your own representation, your utility function would have to value the quotation of the quotation. ``Recipes make poor cake when you grind them up and toss them in the batter.''

Then the complications. The black-hole case is flattened into one outcome, the sister's recovery followed by the end of the Earth, so her life can be valued for itself and still be outweighed. In the flat model, instrumental values have ``vanished entirely''. They come back when the model represents causes. If some state B tends to lead to C however B is reached, you can plan by finding a B that leads to C and an A that leads to B: B has instrumental value because it leads to C, and C ``might be terminally valued - a term in the utility function over the total outcome.''

Instrumental value is only a shortcut for computing plans, to be dropped where its regularity fails: no one gets into a car by ripping off the door with a steam shovel. It is a ``leaky abstraction'', ``as we programmers say''. \nb{The phrase is Joel Spolsky's, from 2002.} Complicate the model too early and you may think instrumental values have a life of their own, so that valuing B commits you to B even without C. Or you may think that a consequentialist maximizing fitness would starve unless it had a terminal value for eating, though no one opens car doors all day for fear of being locked out. With the utility function fixed, instrumental values change only when beliefs of fact change: come to believe penicillin causes pneumonia, and its value drops.

In moral arguments, some disputes are about consequences and some about terminal values. People who disagree about whether banning guns lowers crime agree that crime is bad. In angry arguments this gets lost, and each side decides the other must be sociopaths who really want people killed. ``But I do not think an argument about female circumcision is really a factual argument'' about a shared value. \nb{Charles Stevenson had divided moral disputes into disagreement in belief and disagreement in attitude in 1944. The reasons for circumcision that the World Health Organization reports include both values its opponents reject, such as premarital virginity, and beliefs of fact, such as that the practice has religious support.}

``I fear'' the brain does not keep moral beliefs about means and ends apart: ``We should ban guns'' and ``We should save lives'' do not feel different. My evidence is how they feel. ``Reducing crime'' serves further values, and someone who defends gun rights may also treat ``freedom'' as a terminal value, so even my gun example may involve different terminal values. We cannot print out our network of values, and probably do not store how they got there. Moral dilemmas of the form ``Would you do X if Y'' help, with their own pitfalls. We can reconstruct our values only by ``error-prone projects of cognitive archaeology''. The simple model shows how easy this ought to be in principle, and says nothing of the reward system's further complications, such as the different pleasures of eating chocolate and of anticipating it. ``Being ignorant of your own values may not always be fun, but at least it's not boring.''
